\documentclass[12pt]{article}
\usepackage[T1]{fontenc}
\usepackage[utf8]{inputenc}
\usepackage{lmodern}
\usepackage{textcomp}
\usepackage{enumitem}
\usepackage{geometry}
\geometry{margin=1in}
\begin{document}
\begin{center}
Answer Key - Political Science - Undergraduate Level Assessment 15 \
\small (Answers are ordered exactly as in the question paper.)
\end{center}

\section*{Multiple Choice}
\begin{enumerate}[label=\arabic*.]
\item \textbf{Answer:} C
\\ \textit{Options:} A) To reward recipients for the cessation of negative behaviour, such as human rights abuses.; B) To incentivise recipients to act in ways beneficial to the donor.; C) Compassion for the human suffering of others.; D) To influence recipients or potential recipients through granting or denying aid.
\\ \textit{Note:} The correct answer is "Compassion for the human suffering of others" because economic policies for security are primarily driven by strategic interests and national benefits rather than humanitarian concerns.
\item \textbf{Answer:} D
\\ \textit{Options:} A) distributing propaganda to media; B) endorsing candidates for office; C) socializing with government officials; D) acquiring corporate campaign donations for candidates
\\ \textit{Note:} Acquiring corporate campaign donations for candidates is not a direct method of persuasion by lobbyists, as it involves financial support rather than the more traditional forms of advocacy and influence, such as providing information or mobilizing public opinion.
\item \textbf{Answer:} B
\\ \textit{Options:} A) Assurance of US parity with other powers; B) Striving for American dominance in the international system; C) Putting America's interests first; D) Promotion of America as the world's first democracy
\\ \textit{Note:} The grand strategy of 'primacy' refers to the pursuit of a dominant position for the United States in the international system, emphasizing military superiority and influence to deter rivals and shape global affairs in favor of American interests.
\item \textbf{Answer:} A
\\ \textit{Options:} A) income equality; B) democracy; C) individual liberty; D) due process
\\ \textit{Note:} Income equality is not considered a core value of American political culture, which typically emphasizes individualism, liberty, and equal opportunity rather than the redistribution of wealth.
\item \textbf{Answer:} D
\\ \textit{Options:} A) front-loading; B) prior restraint; C) reapportionment; D) gerrymandering
\\ \textit{Note:} Gerrymandering is the manipulation of electoral district boundaries to favor a particular political party, making it the correct term for the practice described in the question.
\end{enumerate}

\section*{True / False}
\begin{enumerate}[label=\arabic*.]
\setcounter{enumi}{5}
\item \textbf{Answer:} True
\\ \textit{Options:} A) True; B) False
\\ \textit{Note:} The president has the inherent power to declare a state of emergency following a natural disaster as it falls under their constitutional authority to ensure national security and respond to crises effectively.
\item \textbf{Answer:} False
\\ \textit{Options:} A) True; B) False
\\ \textit{Note:} The War Powers Act of 1974 was designed to limit the president's ability to engage in military actions without congressional approval, aiming to reassert legislative authority over war-making powers.
\item \textbf{Answer:} True
\\ \textit{Options:} A) True; B) False
\\ \textit{Note:} Congress has the authority to establish and organize federal courts under Article III of the U.S. Constitution, allowing it to create new courts as necessary to address the judicial needs of the federal court system.
\item \textbf{Answer:} True
\\ \textit{Options:} A) True; B) False
\\ \textit{Note:} The Bill of Rights, particularly through the interpretations of the First, Fourth, and Ninth Amendments, establishes a framework that has been interpreted by the Supreme Court as supporting a right to privacy, which was crucial in the decision of Roe v. Wade regarding a woman's right to make decisions about her own body.
\item \textbf{Answer:} True
\\ \textit{Options:} A) True; B) False
\\ \textit{Note:} The statement is true because the Rules Committee in the House of Representatives is responsible for establishing the procedures and guidelines that govern how legislation is debated, amended, and voted on in the House, thus playing a crucial role in shaping the legislative process.
\end{enumerate}

\section*{Long Form Response}
\begin{enumerate}[label=\arabic*.]
\setcounter{enumi}{10}
\item \textbf{Answer:} Braun and Chyba (2004) argue that indigenous nuclear programs can complicate proliferation safeguards by allowing states to develop the technical capabilities and knowledge necessary for nuclear weapons, potentially undermining global non-proliferation efforts. Their perspective is supported by historical examples where nations with civilian nuclear programs have pursued or enhanced their military capabilities. For instance, countries like South Africa and Iran illustrate how peaceful nuclear ambitions can lead to weapons development. However, this view is challenged by the argument that robust international oversight and cooperation can mitigate these risks, suggesting that the presence of indigenous programs does not inevitably lead to proliferation. Overall, while Braun and Chyba highlight significant concerns, the effectiveness of safeguards and international agreements plays a crucial role in shaping outcomes.
\\ \textit{Note:} correct because it accurately reflects Braun and Chyba's argument that indigenous nuclear programs can lead to the development of technical expertise that may facilitate the pursuit of nuclear weapons, supported by relevant historical examples.
\item \textbf{Answer:} Extraordinary measures are characterized by their deviation from the standard political processes and norms, often arising in response to crises or exceptional circumstances. Unlike ordinary political actions, which operate within established frameworks of law and governance, extraordinary measures may involve unilateral decisions by leaders, suspension of civil liberties, or emergency powers that bypass legislative oversight. These actions can lead to significant implications in the political realm, such as the erosion of democratic norms, increased executive power, and potential abuses of authority. The use of extraordinary measures often reflects the tension between the need for swift action in times of crisis and the principles of accountability and rule of law. Thus, while they may provide immediate solutions, they raise critical questions about the balance of power and the preservation of democratic values.
\\ \textit{Note:} The answer correctly identifies that extraordinary measures diverge from established political norms and processes, often arising in crises and involving unilateral actions, which can have profound implications for governance and civil liberties.
\end{enumerate}

\vfill
\noindent\textit{End of Answer Key}
\end{document}